

Mobile networks are under constant pressure to deliver reliable service, while the resources required to do so, including frequency spectrum, infrastructure, and capacity, remain costly. When operators fail to anticipate future traffic demand, the consequences are wasted capacity, degraded user experience, and inefficient infrastructure investment. Beyond immediate performance, accurate traffic forecasting enables operators to make informed decisions about where and when to upgrade network infrastructure, ensuring that capital is directed toward cells where demand genuinely warrants it \\

%Mobile network cells do not follow a single traffic pattern. Geographic location strongly influences both the level and timing of traffic. These differences are driven by human behavior, such as commuting and recreational activity, as well as network effects like load shifting between nearby cells. As a result, forecasting across all cells is a diverse time-series problem. \\

Mobile network cells do not follow a single traffic pattern. A common assumption in network planning is that geographic location is a strong determinant of seasonal traffic behaviour, for instance, that recreational areas are inherently seasonal while urban areas are not. These differences are driven by human behavior, such as commuting and recreational activity. As a result, forecasting across all cells is a diverse time-series problem, and it remains an open question whether geographic labels alone are a reliable basis for grouping cells by their underlying traffic behaviour. \\

Because of this variation, building a robust forecasting model is not straightforward. One approach is to model each cell independently. This preserves local structure but is computationally expensive and can lead to overfitting, especially for cells with limited historical data. Another approach is to train a single global model across all cells. This allows shared representations of traffic patterns and can capture differences between cell types. However, it reduces interpretability and statistical consistency, since the model combines data from sites with different underlying traffic patterns. A third approach is to pre-classify cells based on their dominant traffic behavior, grouping cells with similar temporal structures. This reduces diversity within each group, making traffic patterns more consistent across cells within the same category. In this setting, simpler statistical models such as ARIMA or OLS could remain competitive when applied to more consistent subsets of cells, compared to more advanced machine learning models. \\


This thesis investigates whether such a data-driven classification can recover, the seasonal grouping that geography is assumed to provide. Specifically, it addresses two research questions: First, can a data-driven statistical method reliably classify mobile network cells into seasonal and non-seasonal categories based on traffic behaviour alone, independent of geographic labels? Second, does training forecasting models on these subgroups yield better predictive performance than a single global model trained across all cells? \\

%This thesis addresses two research questions: First, can a data-driven statistical method reliably classify mobile network cells into seasonal and non-seasonal categories based on traffic behaviour alone? Second, does training forecasting models on these homogeneous subgroups yield better predictive performance than a single global model trained across all cells? \\

The remainder of this work first describes the dataset and preprocessing steps, then introduces a method for classifying cells based on seasonality and traffic behavior, and finally evaluates several forecasting approaches, including OLS, ARIMA, and XGBoost, in terms of forecasting performance. \\